\notechapter{N-20240323}{Combinatorics: Identities, Generating Functions, Catalan Numbers, and More}

These exercises move from binomial identities and recurrences to convexity and graph walks.  What first looks like a mixed problem set gradually becomes a catalogue of ways to encode counting: by subsets, generating functions, recursive decompositions, convex combinations, and matrix powers. The Catalan-number derivations and their combinatorial models follow Alexander Postnikov's MIT 18.212 algebraic-combinatorics lectures.\footnote{\href{https://ocw.mit.edu/courses/18-212-algebraic-combinatorics-spring-2019/pages/lecture-notes/}{MIT 18.212, \emph{Algebraic Combinatorics}}, lecture notes.}

\section{Binomial identities}

Several binomial identities become clearer once their combinatorial meanings are made explicit.  The basic identity
\[
  \sum_{i=0}^n \binom ni=2^n
\]
counts all subsets of an $n$-element set: choosing a subset means choosing its size $i$ and then choosing $i$ elements.

Pascal's identity
\[
  \binom ni+\binom n{i+1}=\binom{n+1}{i+1}
\]
can be proved by fixing a distinguished element in an $(n+1)$-element set.  An $(i+1)$-subset either contains the distinguished element or it does not.  The two cases give the two terms on the left.

Vandermonde's identity
\[
  \sum_i \binom mi \binom n{r-i}=\binom{m+n}{r}
\]
counts choosing $r$ balls from two boxes with $m$ and $n$ balls.  If $i$ balls are chosen from the first box, then $r-i$ are chosen from the second.

Another useful identity is
\[
  \binom nr \binom{n-r}{k-r}=\binom nk \binom kr.
\]
Both sides count a pair $(R,K)$ with $R\subseteq K\subseteq[n]$, $|R|=r$, and $|K|=k$.

\section{Good subsets}

A nonempty subset $A\subseteq\{1,2,\ldots,n\}$ is called good if
\[
  |A|\le \min_{x\in A}x.
\]
Let $a_n$ be the number of good subsets.  The problem asks one to prove
\[
  a_{n+2}=a_{n+1}+a_n+1.
\]

A clean proof splits good subsets of $[n+2]$ into cases.

First, if $n+2\notin A$, then $A$ is just a good subset of $[n+1]$, giving $a_{n+1}$ possibilities.

Second, if $A=\{n+2\}$, this gives one extra subset.

Third, suppose $n+2\in A$ and $|A|\ge2$.  Remove $n+2$ and subtract $1$ from every remaining element.  If
\[
  A'=A\setminus\{n+2\},
\]
then the condition becomes
\[
  |A'|+1\le \min A'.
\]
After shifting $B=\{x-1:x\in A'\}$, this is exactly
\[
  |B|\le \min B,
\]
so $B$ is a good subset of $[n]$.  This gives $a_n$ possibilities.  Hence
\[
  a_{n+2}=a_{n+1}+a_n+1.
\]

A summation proof fixes $|A|=k$.  Every element must then be at least $k$, so the remaining data are chosen from a tail interval, and repeated Pascal identities yield the recurrence.

\section{Generating functions}

For a sequence $(a_n)_{n\ge0}$, its ordinary generating function is
\[
  f(x)=\sum_{n\ge0}a_nx^n.
\]
The note emphasizes a key principle: coefficients can encode counting problems.

For example, the number of nonnegative integer solutions of
\[
  x_1+x_2+\cdots+x_k=n
\]
is the coefficient of $x^n$ in
\[
  (1+x+x^2+\cdots)^k=\frac1{(1-x)^k}.
\]
Using the binomial expansion,
\[
  \frac1{(1-x)^k}=\sum_{n\ge0}\binom{n+k-1}{k-1}x^n,
\]
so the number of solutions is
\[
  \binom{n+k-1}{k-1}.
\]

The notes also use generating functions to separate parity or color constraints.  A typical trick is to multiply several series, each recording the allowed choices for one box or color, and then read off the coefficient of the target power.

% N-20240323 overflow relief A
\enlargethispage{4\baselineskip}
\section{Independent Trials and Generating Functions}

For a nonnegative integer-valued random variable \(X\), its probability generating function is
\[
 G_X(t)=\mathbb E[t^X]=\sum_{k\ge0}\Pr(X=k)t^k.
\]
If \(X,Y\) are independent, then
\[
 G_{X+Y}(t)=\mathbb E[t^{X+Y}]
 =\mathbb E[t^X]\mathbb E[t^Y]=G_X(t)G_Y(t).
\]
For example, one Bernoulli(\(p\)) trial has generating function \(1-p+pt\). Two independent batches of \(2\) and \(3\) such trials have respective generating functions \((1-p+pt)^2\) and \((1-p+pt)^3\). Their sum therefore has generating function
\[
 (1-p+pt)^5,\qquad
 \Pr(X+Y=2)=\binom52p^2(1-p)^3.
\]
This is the probabilistic analogue of ordinary generating-function convolution: independence turns counting or probability of a sum into multiplication of series.

\section{Catalan numbers}

The note then turns to Catalan numbers.  One model is the number of ways to triangulate a convex $(n+2)$-gon:
\[
  C_n=\frac1{n+1}\binom{2n}{n}.
\]
For example, a hexagon has $14$ triangulations.

Let $h_n$ denote the number of triangulations of an $(n+2)$-gon.  Choose the triangle incident with a fixed side.  If its third vertex cuts the polygon into two smaller polygons, the numbers of triangulations multiply.  Therefore
\[
  h_n=\sum_{i=0}^{n-1}h_i h_{n-1-i},\qquad h_0=1.
\]
If
\[
  H(x)=\sum_{n\ge0}h_nx^n,
\]
then the recurrence gives
\[
  H(x)=1+xH(x)^2.
\]
Solving the quadratic and choosing the branch with $H(0)=1$ gives
\[
  H(x)=\frac{1-\sqrt{1-4x}}{2x},\qquad
  h_n=C_n.
\]

This derivation explains why Catalan numbers appear in triangulations, parentheses, Dyck paths, and many noncrossing structures: all of them are recursively split into a left and a right part.

\section{Convex sets and convex hulls}

The geometry part defines a convex set $S$ by the condition that for any $x,y\in S$ and any $t\in[0,1]$,
\[
  tx+(1-t)y\in S.
\]
Equivalently, $S$ contains the line segment between any two of its points.  The note contrasts convex and non-convex examples.

For finitely many points $x_1,\ldots,x_k$, their convex hull is
\[
  \left\{\sum_{i=1}^k t_ix_i:\ t_i\ge0,\ \sum_{i=1}^k t_i=1\right\}.
\]
Induction on $k$ proves the claim: the case $k=2$ is the definition of a segment, and the induction step writes the convex combination as a combination of one point with a convex combination of the remaining points.

\section{Separating a point from a closed convex set}

Let \(C\subset\mathbb R^n\) be nonempty, closed and convex, and let \(x\notin C\). Because squared distance to \(x\) is continuous and grows without bound at infinity, it attains a minimum at some \(y\in C\). For any \(z\in C\), convexity gives \(y+t(z-y)\in C\) for \(0\le t\le1\). The function
\[
 \varphi(t)=\|x-y-t(z-y)\|^2
\]
has its minimum at \(t=0\). Its right derivative therefore satisfies
\[
 0\le \varphi'(0)=-2\langle x-y,z-y\rangle.
\]
Let \(\ell(v)=\langle x-y,v\rangle\). Then every \(z\in C\) obeys
\[
 \ell(z)\le\ell(y)<\ell(x),
\]
since \(\ell(x)-\ell(y)=\|x-y\|^2>0\). Thus an affine hyperplane strictly separates the point \(x\) from \(C\). The proof also explains the geometric role of projection onto a convex set.

For two arbitrary disjoint convex sets, a closest pair need not exist; additional closedness, compactness or other hypotheses must be checked before applying this proof to both sets.

\section{Planar graphs and Euler-type counting}

Graph theory and planar configurations provide another counting mechanism: a planar drawing constrains vertices, edges, and faces through Euler-type formulas.

This connects naturally with the previous convex-geometry material: triangulations, planar graphs, polygon decompositions, and convex hulls are different languages for controlling incidence data.

\section{Graphs and adjacency matrices}

A graph $G$ consists of a vertex set $V=\{v_1,\ldots,v_n\}$ and an edge relation.  For a simple graph, the adjacency matrix $A(G)=(a_{ij})$ is defined by
\[
  a_{ij}=
  \begin{cases}
  1,&\text{if there is an edge from }v_i\text{ to }v_j,\\
  0,&\text{otherwise.}
  \end{cases}
\]
For a multigraph, $a_{ij}$ records the number of edges from $v_i$ to $v_j$.

The source quotes the standard theorem: the $(i,j)$-entry of $A(G)^\ell$ equals the number of walks of length $\ell$ from $v_i$ to $v_j$.  The proof is matrix multiplication.  For $\ell=2$,
\[
  (A^2)_{ij}=\sum_k a_{ik}a_{kj},
\]
which counts all intermediate vertices $v_k$.  The general case follows by induction.

% N-20240323 overflow relief B
\enlargethispage{4\baselineskip}
\section{Catalan objects share a splitting rule}

The parenthesizations of \(n+1\) ordered factors and the triangulations of a convex \((n+2)\)-gon are both counted by \(C_n\). Choosing the final binary operation or the triangle incident with a fixed edge splits the remaining object into independent left and right pieces. Both satisfy
\[
 C_0=1,\qquad
 C_n=\sum_{i=0}^{n-1}C_iC_{n-1-i}.
\]
A separate derivation uses Dyck paths. There are \(\binom{2n}{n}\) paths with \(n\) upward and \(n\) downward steps. The reflection principle bijects those that ever cross below the horizontal axis with paths containing \(n+1\) downward and \(n-1\) upward steps, giving
\[
 C_n=\binom{2n}{n}-\binom{2n}{n+1}
 =\frac{1}{n+1}\binom{2n}{n}.
\]
Recursion, generating functions and reflection now calculate the same number by distinct mechanisms; there is no need to restart the Catalan definition in several chapters of the same note.

\section{Adjacency powers count walks}

For a finite undirected graph with adjacency matrix \(A\), the entry
\[
 (A^k)_{ij}
\]
counts length-\(k\) walks from vertex \(i\) to vertex \(j\). For \(k=2\), this is
\[
 (A^2)_{ij}=\sum_hA_{ih}A_{hj},
\]
which counts intermediate vertices. Inductively, a walk of length \(k+1\) consists of a length-\(k\) walk followed by one edge; matrix multiplication counts exactly these decompositions.

If the graph is \(d\)-regular with \(d>0\), then \(P=A/d\) is the transition matrix for simple random walk. Because \(P\) is symmetric and each row sums to one, the uniform distribution is stationary. If the graph is also connected and non-bipartite, its remaining eigenvalues have modulus strictly less than one:
\[
 1=\lambda_1(P),\qquad
 \rho=\max_{2\le j\le n}|\lambda_j(P)|<1.
\]
Expanding a probability vector in an orthonormal eigenbasis shows that its deviation from the uniform distribution decreases in Euclidean norm by at least the factor \(\rho^k\) after \(k\) steps. In a bipartite connected graph, the eigenvalue \(-1\) causes period-two oscillation, so convergence without averaging fails.

This distinction matters: the graph's second-largest signed eigenvalue is not enough; mixing is controlled by the largest \emph{absolute} nontrivial eigenvalue.

\section{Which encoding should a problem use?}

Subsets and chains naturally invite binomial coefficients; independent decomposition produces products of generating functions; recursively split polygon triangulations produce Catalan numbers; convex hulls encode finite weighted averages; graph walks are encoded by powers of an adjacency matrix. The point of identifying the encoding is that it yields a calculable identity or a proof, with its conditions made explicit, rather than an analogy based on shared vocabulary.
